\begin{Exercise}[label={ex:basso_medio_alto}]
		Scrivere un programma che legga da input un numero intero e stampi su output:
		\begin{itemize}
		  \item la stringa \verb+basso+ se il numero \`e compreso tra $0$ e $3$;
		  \item la stringa \verb+MediO+ se il numero \`e compreso tra $4$ e $8$;
		  \item la stringa \verb+ALTO!+ se il numero \`e compreso tra $9$ e $10$;
		  \item la stringa \verb+Numero non valido+ altrimenti.
		\end{itemize}
		
\end{Exercise}

\begin{Answer}[ref={ex:basso_medio_alto}]
    \lstinputlisting[style=cstyle]{src/basso.c}
    Oppure nella versione pi\`u ottimizzata:
    \lstinputlisting[style=cstyle]{src/basso2.c}
\end{Answer}